% what is deep research 
% frontier lab related work 
\section{Related Work}
\paragraph{Deep Research Agents.}

Recent advancements in agentic reasoning have shifted LLM capabilities from basic tool use toward long-horizon autonomous problem solving, exemplified by deep research. These tasks require agents to perform iterative search, aggregate evidence, and conduct multi-step reasoning, and have emerged as a key frontier for LLMs. Major AI labs have developed proprietary deep research systems, such as OpenAI Deep Research~\citep{openai2025deepresearch}, Claude Research~\citep{anthropic2025claude_research}, Kimi-Researcher~\citep{moonshot2025kimi_researcher}, Grok DeepSearch~\citep{xai2025grok3}, all of which extend LLMs with agentic tool use and long-horizon reasoning tailored for in-depth research. 
Meanwhile, the open-source community has introduced a range of deep research agents. For example, Tongyi DeepResearch~\citep{team2025tongyideepresearch} trains an LLM for long-horizon information seeking via an end-to-end agentic training pipeline, while MiroThinker~\citep{miromind2025mirothinker} explores interaction scaling to support deeper and more frequent agent–environment interactions. Other systems, such as Verl-Tool~\citep{jiang2025verltool}, AgentFlow~\citep{li2025agentflow}, Cognitive Kernel-Pro~\citep{fang2025cognitive}, and DeepMiner~\citep{tang2025beyond}, investigate new training algorithms and dynamic memory mechanisms to enhance research capabilities.
Despite these advances, a critical bottleneck remains: most systems rely on proprietary, live online search APIs, such as Google Search~\citep{google_search_api} and Bing Search~\citep{bing_search_api}). This dependence not only incurs prohibitive costs for large-scale data generation but also introduces reproducibility challenges due to the constantly changing nature of the live web. As a result, scalable and cost-effective synthesis of high-quality research trajectories remains challenging.

% \zhuofeng{TODO: update. consider mention long-horizon trajectory}

\paragraph{Synthetic Environments and Trajectory Generation.}
To mitigate the costs, rate limits, and reproducibility challenges associated with live web interactions, researchers have increasingly turned to synthetic and offline environments for agent training and evaluation. Pioneering works such as WebArena~\citep{zhou2023webarena}, Mind2Web~\citep{deng2023mind2web}, and OSWorld~\citep{xie2024osworld} construct static offline snapshots of web applications and operating systems, providing reproducible testbeds for agentic workflows. In parallel, data synthesis pipelines like AgentTuning~\citep{zeng2024agenttuning} and APIGen-MT~\citep{prabhakar2025apigen} explore generating simulated interaction trajectories to supervise open-weight models. However, these synthetic environments and datasets primarily target short-horizon tool use or simplified web navigation, and thus fall short of the requirements for training \textit{deep research} agents. In particular, they lack the massive, unstructured knowledge corpora and infrastructure needed to simulate dozens of iterative search and reasoning steps. To address this gap, we build an offline synthesis environment around a 15M-document local search corpus. Unlike prior approaches, this environment enables scalable generation of long-horizon deep research trajectories, including many cases that require 100+ tool calls.